% Kappa is 2h(alpha)/(alpha Delta E), not an alpha-independent constant.
\documentclass[11pt]{article}
\usepackage[margin=1.05in]{geometry}
\usepackage{amsmath,amssymb,amsthm}
\usepackage[colorlinks=true,linkcolor=blue,citecolor=blue,urlcolor=blue]{hyperref}

\newtheorem*{admission}{Admission}

\title{\textbf{$\kappa$ Runs with $\alpha$.\\
It Is $2h(\alpha)/(\alpha\Delta E)$.}}
\author{Robert W.\ McGwier\\
\small CTO, Cohere Technology Group\\
\small GitHub: \texttt{n4hy}}
\date{15 August 2026 --- Version 1}

\begin{document}
\maketitle

\begin{abstract}
\noindent
Papers~46--59 used $\alpha=0.02$. There
$S_{\mathrm{global}}=2h(\alpha)$ equals $\delta S$ of
an enclosing ball, and $\kappa=\delta S/P_{\mathrm{flat}}$
looked like one number.

Scan $\alpha\in\{0.005,0.01,0.02,0.04,0.08\}$.
$\kappa$ falls $1.34\to 0.75$ (rel range $0.57$).
It tracks $2h(\alpha)/(\alpha\Delta E)$ at Pearson
$0.999998$. Auditor: $1.26\to 0.85$, C\_lin REFUTED.

At fixed $\alpha$, $\delta S\propto\sum\delta e$ still
holds and the force-law papers still used
$M=P_{\mathrm{flat}}$. Across $\alpha$, $\kappa$ is
not a coupling. The reusable mass is $\sum\delta e$.
\end{abstract}

\begin{admission}
I never scanned $\alpha$. The number I liked was
$2h(\alpha)/(\alpha\Delta E)$ at $\alpha=0.02$. I
did not hide the run.
\end{admission}

\end{document}
